\documentclass[tikz,border=10pt]{standalone}
\input{../../../docs/tikz/dag_preamble.tex}

\begin{document}
\begin{tikzpicture}[
  x=1cm,
  y=1cm,
  every node/.append style={outer sep=4pt, font=\sffamily\scriptsize},
  mwwide/.style={text width=4.45cm, minimum height=1.45cm, align=center, inner sep=3pt}
]

\dagPaperMetadata{Optimal Incentive-Compatible Priority Pricing for the M/M/1 Queue}{Haim Mendelson and Seungjin Whang}{Operations Research 38(5), 870--883}{1990}{https://pubsonline.informs.org/doi/10.1287/opre.38.5.870}
\dagPaperLegendRightOfMetadata{
\node[dag_model, dag_template_legend] (legModel) at (0,0) {source model /\ assumption};
\node[dag_lemma, right=0.35cm of legModel, dag_template_legend] (legDefinition) {source definition /\ formula};
\node[dag_result, right=0.35cm of legDefinition, dag_template_legend] (legResult) {formalized\ result};
}
\daglegend{(legModel)(legDefinition)(legResult)}{Legend}

\begin{dagPaperBody}
\begin{scope}[xshift=0.8cm]

% Economic-pricing column.
\node[dag_model, mwwide] (EconomicModel) at (0,0)
  {\textbf{Assumptions A2, A3, A5}\par value and delay-cost conditions; atomistic individual demand};
\node[dag_lemma, mwwide] (ObjectiveDemand) at (0,-3.2)
  {\textbf{Equations (4)--(5)}\par net-value objective and marginal-user demand equality};
\node[dag_result, mwwide] (TheoremOne) at (0,-6.4)
  {\textbf{Theorem 1}\par optimal price equals the marginal waiting-time externality};

% Stationary queue and homogeneous-pricing column.
\node[dag_model, mwwide] (QueueModel) at (7,-0.0)
  {\textbf{Assumptions A1 and A4}\par stationary Poisson arrivals and a nonpreemptive multiclass M/M/1 queue};
\node[dag_lemma, mwwide] (QueueDefinition) at (7,-3.2)
  {\textbf{Appendix queueing-time definition}\par finite load, residual work, and class priority slacks};
\node[dag_result, mwwide] (AppendixAOne) at (7,-6.4)
  {\textbf{Appendix (A-1)}\par stationary mean queueing-time formula};
\node[dag_lemma, mwwide] (HomogeneousFormulas) at (7,-9.6)
  {\textbf{Equations (8)--(9)}\par homogeneous queueing time, population, and total sojourn time};
\node[dag_lemma, mwwide] (PriorityPrice) at (7,-12.8)
  {\textbf{Equations (10)--(13)}\par priority-dependent price and private expected cost};
\node[dag_result, mwwide] (TheoremTwo) at (7,-16.0)
  {\textbf{Theorem 2}\par strict self-selection at positive stable class flows};

% Heterogeneous PTD-pricing column and its results.
\node[dag_lemma, mwwide] (PTDCriterion) at (14,-6.4)
  {\textbf{Equations (14)--(15)}\par optimal and incentive-compatible PTD equilibrium criterion};
\node[dag_lemma, mwwide] (PTDSchedule) at (14,-9.6)
  {\textbf{Equations (19)--(21)}\par quadratic priority-and-time-dependent price schedule};
\node[dag_lemma, mwwide] (CheatingPenalty) at (14,-12.8)
  {\textbf{Cheating-penalty definition}\par expected total-cost increase from a priority misreport};
\node[dag_result, mwwide] (TheoremThree) at (20,-12.8)
  {\textbf{Theorem 3}\par PTD pricing implements the stable welfare optimum and strict self-selection};
\node[dag_result, mwwide] (TheoremFour) at (20,-16.0)
  {\textbf{Theorem 4}\par zero at the assigned priority and strictly increasing away from it};

\draw[dag_arrow] (EconomicModel) -- (ObjectiveDemand);
\draw[dag_arrow] (ObjectiveDemand) -- (TheoremOne);

\draw[dag_arrow] (QueueModel) -- (QueueDefinition);
\draw[dag_arrow] (QueueDefinition) -- (AppendixAOne);
\draw[dag_arrow] (AppendixAOne) -- (HomogeneousFormulas);
\draw[dag_arrow] (HomogeneousFormulas) -- (PriorityPrice);
\draw[dag_arrow] (PriorityPrice) -- (TheoremTwo);

% Theorem 1 also determines the homogeneous externality price.  The route
% stays in the empty corridor between the economic and queueing columns.
\draw[dag_arrow] (TheoremOne.south) -- ++(0,-1.20) -- ++(3.25,0) |- (PriorityPrice.west);

% The PTD equilibrium criterion uses the same objective and demand semantics.
% Its route crosses only the empty edge corridor between queueing nodes.
\draw[dag_arrow] (ObjectiveDemand.east) -- ++(0.65,0) -- ++(0,-1.55)
  -| (PTDCriterion.west);

% The PTD schedule uses both the externality calculation and the stationary
% heterogeneous priority formula.
\draw[dag_arrow] (TheoremOne.south) -- ++(0,-1.05) -- ++(11.15,0) |- (PTDSchedule.west);
\draw[dag_arrow] (AppendixAOne.south east) -- (PTDSchedule.north west);

\draw[dag_arrow] (PTDSchedule) -- (CheatingPenalty);
\draw[dag_arrow] (PTDCriterion.east) -- ++(1.00,0) |- (TheoremThree.north);
\draw[dag_arrow] (CheatingPenalty) -- (TheoremThree);
\draw[dag_arrow] (TheoremThree) -- (TheoremFour);

\end{scope}
\end{dagPaperBody}
\end{tikzpicture}
\end{document}
